\section{Procedimentos metodológicos}
\label{sec:metodo}
\aescrever
% Modelo: ter, no mínimo, as quatro seções abaixo; dividir conforme a necessidade.
% Fonte de conteúdo: proposta v5, seções 6 e 7; protocolo experimental v5.


\subsection{Classificação metodológica do estudo}
\aescrever


\subsection{População e amostra}
\aescrever


\subsection{Variáveis da pesquisa}
\aescrever


\subsection{Análise e tratamento dos dados}
\aescrever


% O modelo exige, no projeto final (qualificação), o quadro abaixo preenchido.
% Para cada método: onde o texto o apresenta, justifica, valida e testa.
\begin{table}[htbp]
  \caption{Síntese dos procedimentos metodológicos}
  \label{tab:sintese-metodo}
  \begin{tabularx}{\textwidth}{l X X X X}
    \toprule
    \textbf{Método} & \textbf{Apresenta} & \textbf{Justifica} & \textbf{Valida} & \textbf{Testa} \\
    \midrule
    Tipologia                    & & & & \\
    Desenho de pesquisa          & & & & \\
    Unidade de análise           & & & & \\
    Técnicas de coleta de dados  & & & & \\
    Técnicas de análise de dados & & & & \\
    Procedimentos de rigor       & & & & \\
    \bottomrule
  \end{tabularx}
  \fonte{Elaborada pelo autor (2027).}
\end{table}
